\chapter{Conclusiones y Cuestionario}

\section{Respuestas al Cuestionario}

\subsection{Deducción de las Ecuaciones de Lissajous (Ecuaciones I y II)}
Para determinar el desfasaje utilizando el modo X-Y del osciloscopio, se aplican dos señales sinusoidales desfasadas
a los canales X e Y:
\begin{gather}
    x(t) = X_p \sin(\omega t) \\
    y(t) = Y_p \sin(\omega t + \phi)
\end{gather}
Donde $X_p$ e $Y_p$ son las amplitudes máximas de las señales.

Cuando la señal en X cruza por cero ($x(t) = 0$), se debe cumplir que $\omega t = 0$ (o múltiplos de $\pi$). Sustituyendo
este valor de tiempo en la ecuación de la señal Y, se obtiene el punto de intersección con el eje Y, definido
como $A$:
\begin{equation}
    y(0) = Y_p \sin(\phi) \implies A = Y_p \sin\phi
\end{equation}

Por definición, el valor máximo que alcanza la elipse en el eje Y es su amplitud $B = Y_p$. Realizando el cociente
entre ambas magnitudes se obtiene la primera ecuación:
\begin{equation}
    \frac{A}{B} = \frac{Y_p \sin\phi}{Y_p} = \sin\phi \implies \sin\phi = \frac{A}{B}
\end{equation}

A partir de la identidad trigonométrica fundamental $\sin^2\phi + \cos^2\phi = 1$, se deduce la expresión para
el factor de potencia ($\cos\phi$):
\begin{equation}
    \cos\phi = \sqrt{1 - \sin^2\phi} = \sqrt{1 - \left(\frac{A}{B}\right)^2}
\end{equation}
Estas ecuaciones permiten determinar con precisión el desfasaje a partir de la geometría de la elipse.


\subsection{Deducción de la Ecuación del Capacitor de Compensación (Ecuación III)}
La potencia reactiva inductiva de una carga con potencia activa $P$ y desfasaje $\phi$ se expresa como:
\begin{equation}
    Q_L = P \tan\phi
\end{equation}

Para compensar completamente este efecto reactivo, se conecta un capacitor en paralelo cuya potencia reactiva
capacitiva es:
\begin{equation}
    Q_C = V^2 \omega C
\end{equation}
Donde $V$ es la tensión eficaz de la línea y $\omega = 2\pi f$ es la pulsación de la red.

Igualando ambas potencias reactivas para lograr una compensación ideal ($Q_C = Q_L$):
\begin{equation}
    V^2 \omega C = P \tan\phi \implies C = \frac{P \tan\phi}{V^2 \omega}
\end{equation}
Esta expresión corresponde a la fórmula utilizada en el Experimento 2.


\subsection{Deformación de la Onda de Corriente}
En el Experimento 2 se observó una notable deformación de la corriente en comparación con una senoidal pura. Esto
se debe a que la lámpara fluorescente es una carga no lineal. El arco de descarga de gas en el interior del tubo
presenta una resistencia dinámica no lineal (ignición y extinción en cada semiciclo), a lo que se suma el efecto de
saturación magnética del núcleo de hierro del balasto inductivo. Estos fenómenos introducen componentes armónicos de
corriente que deforman la onda fundamental de $50\Hz$.


\subsection{Variación del Factor de Potencia con la Carga Mecánica en Motores}
En el caso de realizar los ensayos sobre un motor de corriente alterna, la aplicación de una carga mecánica en su eje
aumenta el factor de potencia. Esto se debe a que la potencia reactiva de magnetización (necesaria para sostener los
campos magnéticos del estator y rotor) permanece prácticamente constante, mientras que la potencia activa aumenta de
forma proporcional al trabajo mecánico realizado. Dado que el factor de potencia se define como el cociente $P/S$
y $S = \sqrt{P^2 + Q^2}$, un aumento de $P$ con $Q$ constante desplaza el ángulo $\phi$ hacia valores menores, lo
que eleva el $\cos\phi$ acercándolo a la unidad.


\subsection{Cálculo Analítico de las Cotas de Corrección para Ondas Cuadrada y Triangular}
Los voltímetros de respuesta al valor medio están calibrados para indicar el valor eficaz de una senoidal multiplicando
la lectura del valor medio de módulo por el factor de forma sinusoidal ($k_{\text{forma, sen}} = 1,11$).

Para una onda cuadrada bipolar de amplitud $V_p$:
\begin{itemize}
    \item Valor eficaz real: $V_{\text{RMS}} = V_p$.
    \item Valor medio de módulo: $V_{\text{med}} = V_p$.
    \item Indicación del voltímetro de valor medio: $V_{\text{indicado}} = 1,11 \cdot V_{\text{med}} = 1,11 \cdot V_p$.
    \item Error teórico: $e\% = +11\%$.
\end{itemize}

Para una onda triangular bipolar de amplitud $V_p$:
\begin{itemize}
    \item Valor eficaz real: $V_{\text{RMS}} = \frac{V_p}{\sqrt{3}} \approx 0,577 \cdot V_p$.
    \item Valor medio de módulo: $V_{\text{med}} = \frac{V_p}{2} = 0,5 \cdot V_p$.
    \item Indicación del voltímetro de valor medio: $V_{\text{indicado}} = 1,11 \cdot V_{\text{med}} = 0,555 \cdot V_p$.
    \item Error teórico: $e\% = \frac{0,555 - 0,577}{0,577} \cdot 100 \approx -3,8\%$.
\end{itemize}

En la práctica, las mediciones arrojaron diferencias mucho menores ($0,56\%$ y $0,21\%$). Esto se atribuye a que los
multímetros digitales modernos (incluso los de respuesta promedio como el UNI-T UT33A+) emplean técnicas de muestreo
digital y algoritmos internos que atenúan la discrepancia teórica clásica de los rectificadores analógicos pasivos,
además de las limitaciones de ancho de banda y la atenuación de las altas frecuencias de conmutación del generador.


\subsection{Factor de Potencia Menor a 1 en Carga Resistiva Dimmerizada}
En el Experimento 5, al utilizar el dimmer con una lámpara incandescente (carga puramente resistiva), se registró un
factor de potencia de $0,43$. Dado que la carga es resistiva, no existe almacenamiento de energía reactiva en forma
de campos magnéticos o eléctricos (el desfasaje de la frecuencia fundamental es nulo, $\cos\theta_1 = 1$).

Sin embargo, el factor de potencia total se define como el producto del factor de desplazamiento por el factor de
distorsión:
\begin{equation}
    \text{PF} = \frac{P}{S} = \cos\theta_1 \cdot \frac{I_1}{I_{\text{RMS}}}
\end{equation}
Donde $I_1$ es el valor eficaz de la componente fundamental y $I_{\text{RMS}}$ es el valor eficaz total de la corriente.

El dimmer recorta la forma de onda de tensión (y por ende de corriente), lo que genera un contenido armónico sumamente
elevado. Esto incrementa significativamente el valor eficaz total de la corriente ($I_{\text{RMS}}$) respecto de la
componente fundamental ($I_1$). Al reducirse el término de distorsión $I_1/I_{\text{RMS}}$, el factor de potencia cae
drásticamente a $0,43$, a pesar de que la carga es no reactiva.

Cabe aclarar que este bajo factor de potencia no se atribuye al capacitor de la red de retardo RC (usado para el disparo
del Triac). Dicho capacitor suele ser de muy baja capacidad (típicamente del orden de los $100\nF$), lo que a frecuencia
industrial ($50\Hz$) representa una reactancia capacitiva sumamente alta ($X_C \approx 31,8\kohm$). A una tensión de
$220\V$, la corriente reactiva que demanda este circuito de disparo es despreciable (menor a $7\mA$), por lo que su aporte
a la potencia reactiva de desplazamiento ($Q$) no influye en la medición realizada.


\section{Conclusión General}
El desarrollo del trabajo práctico permitió contrastar de forma directa las diferencias operativas entre voltímetros de
respuesta al valor medio y de respuesta True RMS. Se comprobó experimentalmente que, ante la presencia de armónicos y
señales recortadas (como en la salida del dimmer), los instrumentos de valor medio pierden exactitud y subestiman o
sobreestiman los valores eficaces reales, validando la necesidad de emplear equipos True RMS en sistemas electrónicos
modernos. Asimismo, se verificó la efectividad de la corrección del factor de potencia mediante capacitores en paralelo,
reduciendo la corriente de línea en un $62,5\%$ y optimizando la transferencia energética de la red.
